\section{n4. A conditional same-centre obstruction and genuinely wild formal conductor covers}
\textit{Supports: ENTIRE qualified G331; n1.}

First we give a complete conditional obstruction for general conductors.
Let $R$ be an excellent reduced local germ, $S$ its finite normalization
with a single marked NORMAL local centre, and $\mathfrak a$ a primary
gauge. Let $c\ne0$ be a regular conductor element and $b\in S$.
Assume a full conductor-infinite fibre over $o$ has lifts $\alpha,\beta$
with $\alpha(b)=+\infty$, $\beta(b)=A<+\infty$. Assume there is a net
$w_i\to\beta$ on the FULL source section, with $w_i(c)<+\infty$ and
$w_i(b)\to A$. Finally assume a normalized source comparison point $\nu$
and $K>0$ satisfy
\[
 \nu(IS)\le K\,o(I)\quad\text{for EVERY primary/unit }I\subset R.
\]
This is a DOWNSTAIRS-ideal hypothesis, not an unproved estimate for all
$S$ ideals. Put $C=\nu(c)<+\infty$, $B=\nu(b)<+\infty$. For an integer
$L>0$ define
\[
 D_0=1+\min(C+B,L)-\min(C,L),\qquad
 \eta=D_0/K-1-\min(A,L).
\]
If $\eta>0$, an actual continuous signed primary-model obstacle has no
greatest geometric $C$ or $D$ minorant; every continuous minorant and every
canonical degree has uniform error at least $\eta$ from its raw supremum.
The hypotheses explicitly concern a full fibre, ALL supports and ALL
primary ideals. Neither fibre cardinality alone nor non-FG alone implies
this conclusion.

Proof: $cb^n\in R$ for every $n$ because the conductor is an $S$ ideal.
The ideals $(cb)+\mathfrak a^L$ and $(c)+\mathfrak a^L$ are proper
primary (conductor elements are nonunits). Define
\[
 F=-v(\mathfrak a)-v((cb)+\mathfrak a^L)+v((c)+\mathfrak a^L).
\]
For EVERY source lift, write $Y=w(c)$ and $a=w(b)$, allowing infinity.
Then $F=-1-\{\min(Y+a,L)-\min(Y,L)\}$; only finite truncations are
subtracted. It is in $[-1-L,-1]$ and is $-1$ on the conductor-infinite
stratum. On the actual primary blowup of
\[
 T_n=\mathfrak a^n(\mathfrak a^{Ln}+(cb^n))
\]
the negative effective relatively ample divisor $-E_{T_n}/n$ is positive
on EVERY contracted integral curve, BEFORE limits. Its profile is
\[
 q_n=-1-\min(Y/n+a,L)\le F.
\]
The inequality follows from
$0\le\min(Y+a,L)-\min(Y,L)\le\min(a,L)$ when finite, and directly
from the truncations in the infinite cases. At each $w_i$,
$\sup_n q_n=-1-\min(w_i(b),L)$.

The all-primary comparison hypothesis, negative real sums, G331 uniform
comparison and actual decreasing infima give
$g(o)\le g(\rho\nu)/K$ for EVERY geometric $C/D$ function.
For $g\le F$, $g(o)\le-D_0/K$.
Any greatest minorant would dominate all $q_n$, hence have values at
$\rho w_i$ with limiting lower bound $-1-\min(A,L)$, contradicting USC
at $o$. Both raw suprema are equal and finite by n1. Their values on this
net and at $o$ show they are not USC. Every continuous minorant converges
along the net to its value at $o$, giving uniform error at least $\eta$.
Canonical profiles have actual ample positivity by n1 and obey this bound
in EVERY degree. Uniform perturbations $\|F'-F\|<\epsilon$ retain failure
whenever
\[
 (K+1)\epsilon<D_0-K(1+\min(A,L)),
\]
by using $q_n-\epsilon$ and the same comparison at $\nu$.
This proves the conditional theorem, without assuming source CE, local
etaleness or psh positivity of a finite-fibre minimum.

We next verify all its hypotheses for a substantially broader explicit
family than the pure power covers. Fix ANY field $k$, $e\ge2$, $d\ge3$,
$Z=(z_1,\ldots,z_{d-3})$, and
\[
 P(T)\in k[[T]],\quad \operatorname{ord}_T P=e,
 \quad P'(T)\ne0.
\]
Let
\[
 A=k[[u,y,s,Z]],\quad S=k[[t,y,s,Z]],\quad u=P(t),
 \qquad R=A+yS.
\]
The finite local conductor cover is generically SEPARABLE of degree $e$,
including wild ramification when the characteristic divides $e$.
No ground-field roots of unity or perfection are assumed.

Here is a direct proof of the ring and separability assertions. Write
$P=t^eU(t)$ with $U(0)\ne0$. Over $k[[y,s,Z]]$, reduce any series modulo
$P$ to its polynomial of $t$ degree less than $e$, divide the remainder
by $P$, and repeat. The $P$-adic and $t$-adic topologies coincide.
This yields a UNIQUE expression
\[
 S=A\oplus tA\oplus\cdots\oplus t^{e-1}A.
\]
Uniqueness follows by reduction modulo $t^e$, then division by $P$ and
iteration. Thus $A\hookrightarrow S$ is injective finite free of rank $e$.
The submodule $R=A\oplus ytA\oplus\cdots\oplus yt^{e-1}A$ is closed
under multiplication because $yS$ is an ideal of $S$. It is finite over
excellent complete $A$, Noetherian and complete local; modulo the base
maximal ideal its extra generators $yt^j$ are nilpotent, so its only
maximal ideal is
\[
 \mathfrak m=(u,y,s,Z,yt,\ldots,yt^{e-1}),\qquad
 \mathfrak mS=(P(t),y,s,Z).
\]
It is an excellent DOMAIN of dimension $d$. The regular normal ring $S$
is finite over $R$, and $t=(yt)/y$ is in its fraction field. Therefore
$S$ is exactly the finite normalization of $R$, with ONE marked centre.

Its conductor is EXACTLY $yS$. Modulo $y$, the subring inclusion is
$k[[u,s,Z]]\subset k[[t,s,Z]]$, still finite free of rank $e$ and hence
of quotient fraction-field degree $e>1$. If a nonzero element $a$ of the
smaller domain multiplied the larger ring into the smaller, then every
larger element would be $(ab)/a$ in the smaller fraction field, a
contradiction. Any conductor residue modulo $y$ must therefore be zero.
Conversely $yS$ certainly is conductor. Thus
\[
 R/(R:S)=k[[u,s,Z]],\quad S/(R:S)=k[[t,s,Z]],\quad u=P(t).
\]
This proof uses quotient fraction fields and does not replace the
conductor by a primary ideal.

To justify the nontrivial algebraic arc, first work in one variable.
The characteristic polynomial $Q(T,u)$ of multiplication by $t$ on
$k[[t]]$ over $k[[u]]$ is monic of degree $e$, with $Q(T,0)=T^e$.
It is the minimal polynomial over $k((u))$ since the fraction-field
degree is $e$. In $k[[T,u]]$ its vanishing under $u=P(T)$ implies
\[
 Q(T,u)=H(T,u)(u-P(T)),\qquad H(T,0)=-T^e/P(T),
\]
where $H$ is a unit. The kernel assertion follows simply by the formal
change of variable $u'=u-P(T)$, whose inverse is $u=u'+P(T)$.
Differentiating at $u=P(t)$ gives
$Q_T(t,P(t))=-H(t,P(t))P'(t)\ne0$. Thus the degree-$e$ field extension
is separable. Identify one copy of $t$ with $s$ and choose a different
conjugate $\tau$ in a finite normal closure $L$ of $k((s))/k((u))$.
Then $\tau\ne s$.

No unmentioned valued-field completeness theorem is needed here.
G331's quotient-field valuation-extension construction gives an extension
of the $s$-adic valuation to $L$. The integral closure of excellent
$k[[s]]$ in $L$ is finite, normal and one-dimensional. As a finite module
it is $s$-adically complete. Its semilocal residue decomposition lifts
idempotents: Newton iteration for $x^2-x$ squares the error at each step;
$2x-1$ is a unit, also in characteristic two. Completeness gives actual
idempotents. Since the ring is a domain, there is only one maximal ideal.
A one-dimensional normal Noetherian local domain is a DVR, using the same
ordinary normal-curve background as G331. Every element of its fraction
field is a uniformizer power times a unit; the extended valuation is
therefore its order up to the normalization $v(s)=1$. Its fraction field
is complete: a Cauchy sequence is bounded below in value, and after a
fixed uniformizer multiplication lies in this complete DVR. Call its
valuation ring $\mathcal O_L$.

All conjugates of $t$ are integral over $k[[u]]$ and have positive value:
$Q(T,0)=T^e$ forces residue zero when $v(u)>0$. The unit identity above
then gives convergent $P(\tau)=P(s)$, and
$v(P(\tau))=e\,v(\tau)=v(P(s))=e$. Thus
\[
 v(\tau)=1,\qquad a=v(\tau-s)\in[1,+\infty)
\]
is FINITE. Substitution of $t\mapsto s$, $y,Z\mapsto0$ defines a
centred semivaluation $\alpha$ on $S$. Substitution
$t\mapsto\tau$, $y,Z\mapsto0$ into $\mathcal O_L$ defines $\beta$.
Both have gauge one because $s$ has value one. They kill $yS$, and their
restrictions to $R$ agree on EVERY element: both send $u$ to $P(s)$,
$s$ to $s$, $Z$ to zero, and kill $yS$. Their common restriction is $o$.
They are distinct since $\alpha(t-s)=+\infty$ and $\beta(t-s)=a$.
All their nonminimal support primes are retained. No assertion that these
are the only fibre points is needed.

For $r\ge1$, substitute $t\mapsto\tau$, $Z\mapsto0$ into
$\mathcal O_L[[y]]$ and give $y$ Gauss weight $r$. For a nonzero series
$\sum_j c_jy^j$ its value is $\min_j(v(c_j)+jr)$, which is attained:
$c_j\in\mathcal O_L$ and $r>0$. Multiplicativity follows because its
finite initial polynomial lies in the domain
$\operatorname{gr}_v(\mathcal O_L)[Y]$; two initial polynomials have
nonzero product. Thus this is a genuine semivaluation $w_r$ on $S$,
with $w_r(y)=r$, $w_r(s)=1$, $w_r(t-s)=a$, normalized by $\mathfrak mS$.
Every formal substitution converges because $v(t)=v(s)=1$.
For EVERY fixed $f\in S$, its $y^0$ coefficient has finite value $N$ or
is zero. If finite, for $r>N$ all positive-$y$ terms have larger value,
so $w_r(f)=N$; if zero, $w_r(f)\ge r\to+\infty$.
Therefore $w_r\to\beta$ and $w_r|_R\to o$ in the actual ALL-element
boundary topology. This proves the required net (in fact a sequence),
not just convergence on generators.

Choose an integer $L>1+a$. In the regular formal coordinates
$(s,b=t-s,y,Z)$, let $\nu$ be the generic monomial valuation with
weights $(1,L,1,1,\ldots,1)$. It has gauge one since $\nu(P(t))=e$
and $\nu(y)=\nu(s)=1$. For EVERY downstairs element $f$ with finite
$o(f)=N$, expansion in these coordinates has a nonzero pure $s^N$
monomial: setting $b=y=Z=0$ is exactly the defining evaluation for $o$.
The positive-weight generic monomial valuation is at most $N$; different
monomials do not cancel. A finite generator attaining a primary ideal's
finite $o$ value therefore gives
\[
 \nu(IS)\le o(I)\quad\text{for EVERY primary/unit }I\subset R.
\]
Units have value zero. This verifies the conditional hypothesis with
$K=1$ for ALL ideals, without using a normal-centre Izumi constant or
assuming an estimate for ALL upstairs ideals.

Take $c=y$, $b=t-s$. The explicit obstacle and competitors are
\[
 F=-v(\mathfrak m)-v((y(t-s))+\mathfrak m^L)
       +v((y)+\mathfrak m^L),
 \qquad T_n=\mathfrak m^n(\mathfrak m^{Ln}+(y(t-s)^n)).
\]
They are ACTUAL downstairs ideals, since $yS\subset R$.
On EVERY source lift put $Y=w(y)\ge1$, $b_w=w(t-s)\ge0$.
Then $F=-1-\{\min(Y+b_w,L)-\min(Y,L)\}\in[-L,-1]$.
The upper and lower bounds include all infinite-value supports; if
$Y\ge L$ the bracket is zero, and otherwise it is at most $L-Y\le L-1$.
At $\nu$, $F=-L$, whereas at $o$, $F=-1$.
The all-ideal comparison gives $g(o)\le g(\rho\nu)$ for the ENTIRE
geometric $C/D$ class, so every minorant has value at most $-L$ at $o$.
The constant $-L$ is itself a minorant, hence the raw value is EXACTLY
$p(o)=-L$. Since $A_m$ contains $\mathfrak m^{mL}$ and $h_m$ is a minorant,
$h_m(o)=-L$ in EVERY degree as well. On the approaching points,
\[
 \sup_n q_n(\rho w_r)=-1-a,qquad p(\rho w_r)\ge-1-a.
\]
The conditional theorem consequently proves no greatest $C/D$ minorant,
non-USC raw envelope, and
\[
 \|p-h_m\|_\infty\ge L-1-a>0\quad\text{for EVERY }m\ge1.
\]
The same bound holds for EVERY continuous minorant, and failure is stable
under $2\epsilon<L-1-a$. No exact formula for $p$ away from these points
is asserted. Every $q_n,h_m$ has actual negative effective ample primary
positivity on EVERY contracted curve before limits. No source CE is used.

In particular, for ANY field of characteristic $p>0$ take
\[
 P(t)=t^p+t^{p+1}.
\]
The order is $e=p$ but the derivative is $t^p\ne0$: this is a genuinely
wild separable degree-$p$ cover, outside the previous pure-power radicial
case. With $b=\tau-s\ne0$, the identity $P(s+b)=P(s)$ reads
\[
 b\{b^{p-1}(1+s+b)+s^p\}=0.
\]
Since $1+s+b$ is a unit, $(p-1)a=p$, so $a=p/(p-1)$.
Choosing $L=4$ gives the explicit, all-field bound
\[
 p(o)=-4,qquad
 \|p-h_m\|_\infty\ge 3-p/(p-1)\ge1\quad(m\ge1).
\]
The formal normalization has only one regular local centre, the conductor
is exactly $yS$, and the failure is nonnormal in every dimension $d\ge3$.
Pure powers from 425 are not claimed as new. The genuinely new domain
here includes separable higher-order perturbations whose ramification
index is a pure characteristic power.

For another primary gauge the canonical ideals are unchanged for the
transported obstacle and the gap along this convergent net becomes
$(L-1-a)/v_o(\mathfrak b)$, or $(3-p/(p-1))/v_o(\mathfrak b)$ in the
wild example. This follows by the same positive continuous denominator
at EVERY support; no conductor-value uniform bound is invoked.

This is a complete conditional sufficient obstruction, fully verified on
the stated formal separable-cover family. It is not a necessary/sufficient
classification of every same-centre conductor cover. Normal higher-dimensional
CE is not disproved, source CE in positive characteristic is not presumed,
and literal radial, product, energy and full operator targets stay separate.
Artifacts: notes/n4.tex, P12.3-quotients-and-wild-conductors.tex,
exact-check-input.json and reports/exact-checks.json. Rerun:
python exact\_checks.py; python build\_delivery.py; python verify\_delivery.py.
Finite checks are diagnostics; the all-support assertions are the proof above.
